\section{The integral positive determinant}\label{sec:positive-determinant}

We construct the integral determinant used in Proposition~\ref{prop:positive}.
The construction has three parts: a residual test controls the divisor \((2)\),
finite--singular switches control the other divisors away from the fixed local
factors, and character evaluations remove the remaining possible poles.
The central specialization is computed here in analytic rank zero;
Section~\ref{sec:positive-completion} treats rank one and a center of initially
unknown analytic rank. Throughout the construction itself, no analytic-rank
hypothesis is imposed.

\subsection{Parameters and the determinant statement}

We work with the non-CM curve \(E'\), the square classes \eqref{eq:K1} with \(A=1\), and the integral symbol combination constructed in Section~\ref{sec:siegel}. Optionally we compare simultaneously with the twist \(E'^{(h)}\), \(h>0\) fundamental prime to \(2N\). 

Fix a finite support \(S\) containing \(\infty\), \(2\), the primes dividing \(N\), and the support of the fixed twist when one is used. Write \(S_f\) for its finite part. Take fixed \(c,d>1\) prime to \(6\prod_{q\in S_f}q\), congruent to 1 modulo the necessary fixed orders (\(N\) and also \(h\) if used). They can be avoided at all varying primes below. Use square orders \(M\) whose support is exactly the indicated finite support, including auxiliary places when required, and which are divisible by the requisite fixed orders. Their \(2\)-parts form a genuine tower. 

Throughout \(T\) in a branch denotes the self-dual Tate lattice for the curve being tested (\(E'\) or \(E'^{(h)}\)); geometrically the lattices identify with the fixed quadratic scalar inserted as in \eqref{eq:K4}. Write \(V=T\otimes\mathbb Q_2,\ W=T/2T\).

At stages \(i\) use a bounded tuple \(R_i=(r_{1,i},\ldots,r_{k,i})\) of new distinct odd good primes. Ultimately \(k\ge1\) is fixed on the limit and \(r_{j,i}\to1\) dyadically. Use an even surjective \(2\)-power cyclic quotient of conductor \(r_{j,i}\) for each variable there, orders tending to arbitrarily high depth, and include all these primes in the allowed and omitted support. With the true real cyclotomic \(\mathbb Z_2\)-direction we use the limit ring
\[
                       B=\mathbb Z_2[[t,u_1,\ldots,u_k]],
\]
the scalar \(\chi\) on chosen generators being \(1+t,1+u_j\) respectively. The scalars are first defined at each finite stage, with the finite-order relations on the tame generators, and then passed to Artin matrix limits as in Section~\ref{sec:foundations}. We retain models also over the stage rings with actual \(t\)-variable before those limits. Let \(M_\chi=T\chi\) denote the coefficient and \(z\) the smoothed Kato class via the group-coefficient map, on the full global problem at support \(S\cup R_i\). Orient all evaluations as in weighted corestriction; we allow the consistent inverse variable on the period-law side depending on conventions. 

In simultaneous comparisons include a fixed order-two group-ring variable first for \(\chi_h\) and apply its two evaluations to common diagrams (so in the twist branch the fixed factor then goes into \(T\)). Both evaluations are taken over the finite local group order; no division by two is used to project onto a character.

Use the limiting positive complex
\[
 C^+=\mathrm{fib}\big(C_{\rm glob}(M_\chi)\to C_{\mathbb R}(M_\chi)\big)
\]
with ordinary real cochains. It is perfect by the finite-model construction of Section~\ref{sec:foundations}. The class \(z\) lifts integrally to positive \(H^1\) by the imaginary level observation of \eqref{eq:K1}, also using compactness to take compatible tower lifts. The connecting image of the real basis \(a\) (transported as given even for the positive twist) is integral there. These operations can precede the two order-two evaluations when used.

\begin{proposition}[Integral positive determinant]\label{prop:positive-determinant:integral}
The auxiliary tuples \(R_i\) can be chosen with the following properties.
The construction permits one prescribed sequence among their entries, provided
that it satisfies the varying-support and cyclic-quotient conditions above,
splits on \(W\) and in the fixed twisting field, and has a limiting Tate
Frobenius matrix of determinant one and trace different from two.

The complex \(C^+\otimes_B B_{(2)}\) is a free module of rank two in degree
one, up to homotopy. Generically over \(B\), the tensor \(a\wedge z\) is
nonzero and gives a coordinate \(U_{\rm raw}\in\operatorname{Frac}B\) relative
to a generator of \(\mathcal D(C^+)\). In a simultaneous twist comparison the
generator is inherited from the common diagram before the two character
evaluations. The tensor is independent of the positive lift of \(z\) after
passing to the fraction field.

Define the normalized coordinate by
\[
 U=\frac{U_{\rm raw}}{\Delta}
                    \prod_{\substack{q\in S_f\\q\ne2}}\frac{D_q}{P_q},
 \qquad
 D_q=\det(1-\chi({\rm Fr}_q)\rho_T({\rm Fr}_q)/q\mid V_{I_q}).   \tag{L1}\label{eq:L1}
\]
Here \(I_q\) indicates inertia and the subscript in \(V_{I_q}\) takes coinvariants. \(P_q\) is the \emph{actual} local inverse \(L\)-polynomial at 1 of the curve being tested with the unramified scalar twist in the period-law orientation. \(\Delta\) is the smoothing factor of \eqref{eq:K4} with the actual field weight including the fixed character if present. All exponents in this formula are the limiting exponents. There is no
\(D/P\) correction for a moving place. The assertion is that \(U\in B\setminus\{0\}\).
If the construction is made simultaneously for \(E'\) and a positive
fundamental twist \(E'^{(h)}\), the two reductions of \(U\) differ by a unit
of \(B/2B\). In particular, they are simultaneously units at the closed point.
\end{proposition}

We prove the proposition through the next three subsections. The distinction
between the divisor \((2)\) and the horizontal divisors will allow the latter
argument to use fixed clearing powers of two without losing the former.

\begin{lemma}[The fixed factors]\label{lem:positive-determinant:factors}
The factors \(\Delta,D_q,P_q\) in \eqref{eq:L1} are integral and remain
nonzero modulo two after all tame parameters are set to zero. Each
\(D_q/P_q\) has central value one and reduces to a group-like unit in the
residual fraction field.
\end{lemma}
\begin{proof}

All indicated fixed polynomials and smoothing factors are integral, nonzero modulo 2 even upon \(u_1=\cdots=u_k=0\). Indeed each of the real cyclotomic exponents of \(q,c,d\) here is nonzero, and \(D_q,P_q\) use the same type of elliptic polynomial in possibly opposite orientations (good quadratic, signed multiplicative linear, or trivial additive). This follows from the ordinary local Tate description: unramified action at good primes, quotient coinvariant unramified sign line at multiplicative primes, no rational coinvariant at additive primes. It uses the already twisted curve if testing that branch. Palindromicity gives the asserted residual unit, and the two orientations
agree at the central specialization, giving ratio one.
\end{proof}

\subsection{Auxiliary primes and residual concentration}\label{subsec:positive-determinant:residual}

\begin{lemma}[Residual concentration]\label{lem:positive-determinant:residual}
The auxiliary primes in Proposition~\ref{prop:positive-determinant:integral}
can be chosen so that the positive complex is free of rank two in degree one
over \(B_{(2)}\), up to homotopy. The two raw determinant tensors in a
simultaneous positive-twist comparison have identical reductions in the
residual fraction field.
\end{lemma}
\begin{proof}
\emph{Required local tests.}

All added primes split on \(W\) and in the fixed twisting field, and their arithmetic Tate Frobenius matrices have ultrafilter limits \(g_j\) with determinant 1 and trace \(\ne2\) (use chosen local embeddings). This includes the required hypothesis for the distinguished prime. 

Use Frobenius lifts there and at the other moving primes trivial in their \emph{own} tame cyclic field, and inertia generators taking the chosen generator values. Write \(e_{jl}\in\mathbb F_2\) for the bit of the Legendre symbol \((r_{j,i}/r_{l,i})\) for \(j\ne l\) (stagewise), symmetric by reciprocity. We arrange the following two properties.
\begin{enumerate}[label=(\roman*)]
\item For each \(q\in S_f\), the limiting unramified tame-scalar exponent vector at \(q\) is nonzero.
\item If \(W\) is irreducible, at the stages on the filter, for common generic \(\lambda_j\in\overline{\mathbb F}_2^\times\), the following residual \(W\otimes\overline{\mathbb F}_2\) Selmer problem is zero: unramified at fixed finite and other nonmoving places, unrestricted at infinity, and at the moving places
\[
             f_j=b_j s_j,\qquad b_j=\lambda_j^{-1}\sum_{l\ne j}e_{jl}\lambda_l ,
             \qquad f_j=x({\rm Fr}_j),\quad s_j=x(\sigma_j),       \tag{L2}\label{eq:L2}
\]
with \(\sigma_j\) the indicated tame generator. If \(W\) is reducible use just one dummy prime in addition to the distinguished prime if used; when there are two, require their mutual bit to be one.
\end{enumerate}

\emph{The dummy prime and the fixed-place exponents.}
First add the dummy prime. To achieve (i), prescribe over the field of \(W,\chi_h\) (where used) and the \(2\)-cyclotomic tower simultaneous nonzero Kummer exponents on compatible \(2\)-power roots of each \(q\in S_f\), using accuracy tending to infinity. Each coordinate can indeed be nonzero: otherwise the pure cyclotomic tower would have only bounded-degree further radical extensions for that \(q\), hence none (the bounded image on the compatible roots there is a finite subgroup of \(\mathbb Z_2(1)\)). Even all positive radicals of \(q\) cannot lie in the abelian tower, by Eisenstein and nonreal conjugates. 

Joint nonzero values now follow by additivity, and the bit with the distinguished prime remains free by its separate new ramification (take bit zero when \(W\) irreducible). In tame residue characters at the chosen primes the power residue exponents thus have bounded valuations (normalization of a cyclic generator only changes by units). 

Trace nonexceptionality can simultaneously be kept with \(v_2(2-\operatorname{tr} g_j)\) bounded: within the kernel of such bounded-derived-length constraints the Tate image contains a fixed deep \({\rm SL}_2\) by open image~\cite{Serre72} and closed iterated commutators as in the evaluation construction of Section~\ref{sec:foundations}. Its perturbations can keep the trace off 2 uniformly for a compact set of starting matrices. Chebotarev now gives stage choices to increasing precision, split also to the needed cyclotomic depths.

\emph{Killing the irreducible residual Selmer space.}
We next arrange (ii) when \(W\) is irreducible. Work first over \(\mathbb F_2\), initially and at new primes with all mutual bits zero, so \(f_j=0\). Write \(F,F^*\) for the resulting Selmer and exact-annihilator Selmer subspaces in global \(H^1(W)\), using the Weil pairing. At fixed odd places unramified is self-annihilator by isotropy, dimension \(h^0\), duality and the local Euler formula. At 2 unramified is isotropic of dimension \(h^0\) with annihilator two dimensions bigger. At the moving split primes the planes \(f,s\) pair by the mixed Weil pairing up to unit, separately isotropic (scalar Hilbert pairing alternating since \(r_j\equiv1\bmod4\)). At infinity the dual condition is strict. Thus by Wiles's dimension formula (Poitou--Tate, subtracting ordinary \(h^0\) at each local place)
\[
\begin{gathered}
          h^*=h+\delta,\qquad h=\dim F,\qquad h^*=\dim F^*,\\[3pt]
 \delta=\begin{cases}
 1&\text{if complex conjugation acts nontrivially on }W,\\
 0&\text{otherwise}.
 \end{cases}
\end{gathered}
\]
Indeed in the first case real \(h^0=1,h^1=0\); in the second both equal 2. At a new prime with zero bits, switching there from unramified to \(f=0\) gives
\[
             h_{\rm new}\le h-\operatorname{rank}{\rm loc}_r F
                              +2-\operatorname{rank}{\rm loc}_r F^* .
\]
The strict kernel is unchanged by inflation and the new singular coordinate must annihilate old dual evaluations. This bound works also with changed fixed conditions and their exact duals.

At a new split prime evaluations can be prescribed arbitrarily on a joint \({\rm End}_{G_\mathbb Q} W\)-independent set of residual classes. Indeed restriction to the joint kernel of \(W\) and the pro-2 abelian data for bits, cyclotomic and fixed-twist splitting detects these classes: the quotient has a normal subgroup of order three with no invariants. Additive joint evaluations there form a submodule of a sum of the simple \(W\), hence surject by independence. The bits with old primes are freely imposed by disjoint quadratic ramification over the fixed split fields. Evaluation fields do not upset the trace-limit freedom by the same bounded-derived-length argument. Take Chebotarev representatives as above. Dimensions are uniformly bounded for a bounded tuple, so choices and a bounded iteration based initially on a bounded dimension can be fixed in pattern along the filter.

For cyclic cubic image the spaces with the stated zero-bit conditions are \(\mathbb F_4\)-spaces via that endomorphism field (annihilators also stable); if \(h>0\) arrange rank two over \(\mathbb F_2\) on both. 

For \(S_3\) image arrange ranks \(\min(2,h),\min(2,h^*)\) simultaneously by choosing a basis including the intersection. This drops \(h\) until zero except possibly \(h=h^*=1\). 

If the lines then coincide, the intersection problem (unramified at 2 and strict at infinity) has dual of dimension three (by the formula) containing that line. Switch at a new prime evaluating nontrivially on the line and surjectively on this larger dual. The same bound kills the intersection and keeps \(h=h^*\le1\). 

Once the two lines are independent, add a zero-bit prime \(A'\) evaluating the generators on nonparallel vectors \(v,w\), respectively. If after switching \(h>0\) persists, both spaces still have dimension one and the singular coordinates at \(A'\) of their new respective generators are \(w,v\), by vanishing of strict kernels and orthogonality. 

Add now a last prime \(B'\) evaluating trivially on both new lines but with bit 1 with \(A'\), others zero. With \(x=\lambda_{B'}/\lambda_{A'}\) the conditions are \(f_{A'}=x s_{A'},\ s_{B'}=x f_{B'}\), recovering the preceding problem at \(x=0\). Its generator cannot lift modulo \(x^2\): pair the linear coefficient globally with the preceding dual generator, giving \(\langle w,v\rangle\ne0\) at \(A'\) and zero at every other place. The modified local equations form a matrix over \(k'[[x]]\) on
these fixed global and local spaces. A nonzero generic solution could be
scaled to an integral solution with nonzero reduction. That reduction
would be a nonzero multiple of the preceding one-dimensional Selmer
generator, and rescaling would give the forbidden lift modulo \(x^2\). Thus we have vanishing on a nonempty open as required. 

All avoidance polynomials can be fixed along the filter by bounded dimensions and finite base field; we work with \(\lambda_j\) in a sufficiently large common finite \(k'/\mathbb F_2\).

\emph{A one-parameter residual test.}
Set $A=k'[[z_0]]$ and specialize
\[
  2=t=0,\qquad u_j=\lambda_jz_0.
\]
In the irreducible case choose $(\lambda_j)$ in the nonempty open set
furnished by (ii); in the reducible case retain the one- or two-prime
configuration specified there. Choose the vector also so that each fixed
finite place in (i) has nonconstant unramified scalar. The latter conditions are nonempty: the first
nonzero term of the scalar is determined by the least valuation of its
nonzero exponent vector, and cancellation excludes a proper closed
condition on the $\lambda_j$.

Over $k'((z_0))$, Frobenius at a fixed finite place consequently has
no invariants, on either coefficient. At a moving place the nontrivial
inertia scalar gives the same conclusion. Local duality also gives
$H^2=0$ at every finite place. These are statements about the local
limit models: at fixed places the coefficient actions stabilize at
every precision; at moving places the inertia/residue models give the
assertions. Evaluation on these local groups also shows that global
$H^0$ vanishes.

\emph{The remaining global obstruction.}
It remains to kill the kernel of finite localization on ordinary
degree-one cohomology of the dual coefficient. Here is the precise
Poitou--Tate reduction. At a finite stage and scalar precision $A/(z_0^n)$, write $K_n$ for
that dual strict-at-finite kernel. Subscripts $n$ below denote the
corresponding finite-precision complexes. Finite Poitou--Tate gives
\[
 \dim_{k'}H^2(C_n^+)
 \leq \sum_{v\in S_f\cup R_i}\dim_{k'}H^2(L_{v,n})
       +\dim_{k'}K_n,
\]
and bounds $\dim_{k'}H^3(C_n^+)$ by the dual global invariant
dimension; higher positive cohomology vanishes. To see the first
bound, the cokernel of ordinary $H^1$-restriction to the real place
is dual to the real image of $K_n$. The kernel of $H^2$-restriction
to the real place maps to the finite local $H^2$'s, with kernel dual
to the everywhere-locally-zero subspace of $K_n$. These two
contributions together have dimension at most $\dim K_n$.
Ordinary and Tate-modified real terms agree in the positive degrees
used here. A perfect coefficient-ring functional identifies the
finite duals in this calculation.

Apply this inequality stagewise and then to the limit matrices.
For a finitely generated $A$-module, the dimension of its reduction modulo
$z_0^n$ is its free rank times $n$, up to a bounded term. The local
vanishings just proved therefore remove the local contributions
after division by $n$. The strict-finite kernel is likewise measured
by its mapping-fiber model; finite local $H^0$ has rank zero.
Consequently a rank-zero dual strict-finite kernel implies that
$C^+\otimes_Ak'((z_0))$ is concentrated in degree one.

\emph{Reduction of the dual kernel.}
Let $N$ be the full saturated preimage of this kernel in ordinary
dual $H^1$ over $A$, including its torsion. Saturation and the
cohomology base-change sequence inject $N/z_0N$ into residual
ordinary $H^1$. At each fixed finite place its image is unramified:
inertia acts constantly, so a restriction that becomes a coboundary
after multiplication by a power of $z_0$ has residual coefficient a
coboundary on inertia.

At a moving place the tame cocycle relation on the limit is
\[
 (\chi^{-1}(\sigma_j)-1)x({\rm Fr}_j)
   =(\chi^{-1}({\rm Fr}_j)-1)x(\sigma_j).
\]
Indeed, discard prime-to-two inertia and use $r_j\to1$ in the
Frobenius conjugation relation, including its cocycle sum. Dividing
by $z_0$ and reducing gives \eqref{eq:L2}. These conditions are retained in
the finite diagrams: at a fixed place finitely many inertia
evaluations detect a coboundary, and at a moving place the two
displayed evaluations were marked. Thus the residual tests apply
to the reduction of $N$, not merely to unrelated stage classes.
In the irreducible case (ii) gives $N/z_0N=0$, hence $N=0$.

\emph{The reducible case.}
For a trivial constituent, residual classes unramified away from
the moving finite places are spanned by $[r_j]$, since
$r_j\equiv1\pmod8$. Their evaluations are the Legendre bits in
\eqref{eq:L2}, with own-prime Frobenius bit zero. With one prime, or with
two primes of mutual bit one, the solution space has dimension
one. The saturated preimage $N$ already contains the torsion
cocycle $(\chi^{-1}-1)/z_0$, with nonzero reduction; it can be
defined compatibly at one higher precision. Since this torsion
class occupies the sole allowed residual dimension, $N$ has
free rank zero. The preceding Poitou--Tate argument gives
positive concentration for a trivial constituent. Exact
triangles for extensions then give concentration for $W$.

\emph{Concentration and twist comparison.}
Matrix cancellation now gives concentration over the full residual
fraction field, and then over $B_{(2)}$. The Euler characteristic
of $C^+$ is $-2$: evaluate a stage at constant characteristic-zero
coefficients, where the ordinary global Euler characteristic is
$-1$ and the real invariant line contributes one more subtraction.
Thus the minimal model over $B_{(2)}$ is a free rank-two module
in degree one. The integral tensor $a\wedge z$ has nonnegative
valuation there; the fixed factors are units at $(2)$.

In the simultaneous twist construction the common diagram,
determinant basis, and classes give identical raw reductions in
the residual fraction field. The smoothing reductions agree too.
After the remaining height-one tests prove $U\in B$, the
group-like residual factors $D_q/P_q$ therefore show that the
two reductions of $U$ differ by a unit of $B/2B$.

\end{proof}

At this point integrality has been proved only at \((2)\). We next prove
generic nonvanishing and control all horizontal divisors; the closed-point
unit comparison in Proposition~\ref{prop:positive-determinant:integral} will
then follow from the residual comparison just obtained.


\subsection{Horizontal line switches}\label{subsec:positive-determinant:horizontal}

We use three characteristic-zero tests. Each has its own diagram and
its own old local conditions:
\begin{center}
\small
\begin{tabular}{@{}p{.16\textwidth}p{.23\textwidth}p{.28\textwidth}p{.23\textwidth}@{}}
\toprule
Source & Test ring & Old finite conditions & Use \\
\midrule
(i) Moving \(B\)-limit & \(B_{\mathfrak p}\) & Full at every finite place &
Divisibility away from fixed factors \\[3pt]
(ii) Fixed support, no \(R_i\) & \(\operatorname{Frac}\mathbb Z_2[[t]]\) &
Full at every finite place & Generic nonvanishing \\[3pt]
(iii) Fixed stage, character \(\eta\) &
\((\Lambda_\eta)_{\mathfrak q}\) and its fraction field &
Unramified at odd places, full at two & Removal of fixed divisors \\
\bottomrule
\end{tabular}
\end{center}
In (i), \(\mathfrak p\) is a height-one prime satisfying
\(\mathfrak p\nmid2\Delta\prod_{q\in S_f,\,q\ne2}D_qP_q\).
In (iii), take one sufficiently deep original stage and specialize every
variable at \(r_{j,i}\) to a nontrivial finite character; we call these
characters active. Their product is \(\eta\), and
\(\Lambda_\eta=\mathbb Z_2[\eta][[t]]\), with the true cyclotomic parameter
retained. The test prime \(\mathfrak q\) is horizontal of height one. Additional
restrictions on \(\eta\), and the square or rectangular class used in
that test, are specified below.

The fixed positive quadratic scalar, if active, goes into \(T\). Write \(\mathscr O\) for a test DVR here, or use just its fraction field for a generic-only test; \(2\) is invertible. Use the \textbf{ordinary} global problem (no real condition after inverting 2). In determinant valuation for a positive full complex concentrated as asserted, this simply removes the invariant real basis term: the ordinary real complex after base change is the free invariant line in degree zero, with basis \(a\), since all these characters are even. Thus the coordinate of \(a\wedge z\) has the same valuation as the ordinary determinant coordinate using \(z\). All dualities and inflation comparisons here are after the model limit and \(2\)-inversion as in Section~\ref{sec:foundations} (one can use the comparisons via \(\mathbb Q(i)\)). For a fixed-stage or fixed-support true tower, later switches use a new sequence/ultrafilter internal to the test, keeping its original diagram before switches; true-tower comparisons at the base are then unchanged.

\begin{lemma}[Horizontal switching]\label{lem:positive-determinant:horizontal}
In any of the three test settings, let \(z_{\rm sys}\) be the indicated
square or rectangular class, satisfying the old local conditions. If its
generic image is nonzero, the ordinary complex has a single generic
cohomology line, in degree one. At every indicated horizontal test DVR the
determinant coordinate of \(z_{\rm sys}\) has nonnegative valuation.
\end{lemma}
\begin{proof}
\emph{Old local conditions and duality.}
We first check cohomology at the old places, over both the generic and
residue fields when a DVR is being tested.
\begin{itemize}
\item In the first test dual local \(H^0\) vanishes at every finite place of the old support. At fixed odd \(q\) this follows from \(D_q\) (duality of inertia-invariants with coinvariants, including the Tate twist). 

At 2 choose an inertia element mapping to a generator in the real cyclotomic direction, and a Frobenius lift trivial in that direction. The invariance equations then require the simultaneous vanishing of a nonzero polynomial up to a unit in \(t\) alone (monic after writing the determinant condition with the inverse scalar cleared), and a nonzero series independent of \(t\). The latter uses property (i) of the tame exponents in the residual construction and constant Tate action of the chosen element. The first is regular in \(t\), so its distinguished polynomial has no factor independent of \(t\). It therefore has no height-one prime factor in common with the second series, by Weierstrass division. 

At a moving \(r_{j,i}\) invariance on inertia requires \(u_j=0\); on that divisor the Frobenius determinant condition is not identically satisfied, by \(\det(g_j)=1,\operatorname{tr}g_j\ne2\) on setting scalars centrally. These calculations test the actual local models (fixed groups by coefficient stabilization, moving odd primes by the inertia/residue complexes or Frobenius and inertia evaluations). 

In the second, generic \(t\)-test the same dual vanishing at fixed places uses inertia at 2 and nonconstant \(t\)-Frobenius scalar at every odd \(q\).
\item In the third test we require the active characters such that already a fixed Frobenius lift at 2 trivial in the cyclotomic direction kills dual invariants, by its action with the constant part of the scalar. 

At each old odd \(q\) of the fixed-test support (including all now-fixed tame variable places), put the inflated residue cochains on the inertia invariant block for the unramified complex condition, similarly dually. This is well-defined here after inverting 2 with base change to both fields. Indeed inertia action of \(V\otimes\eta\) is constant; after exact prime-to-2 kernel invariants on inertia, the two-term procyclic \(2\)-inertia linear algebra is constant, split in cohomology over the rational coefficient space before imposing residue Frobenius. The quotient relative to the unramified condition computes residue cochains on the inertia coinvariants with negative Tate twist in degrees \(1,2\), by this calculation and inertia truncation (Frobenius conjugation on the generator evaluation modulo coinvariant relations inserts \(q^{-1}\)). The same formulas hold on the local models in subsequent switches. Its determinant of differential up to unit is
\[
 D_{q,\eta}=\det(1-\chi_t({\rm Fr}_q)\rho_{V\otimes\eta}({\rm Fr}_q)/q
                                 \mid (V\otimes\eta)_{I_q}),
\]
where \(\chi_t\) denotes the \(t\)-scalar. This is nonzero by nonconstant cyclotomic action (and equals 1 at active tame variable conductors, by good reduction there). The quotient has zero \(H^1\) over the DVR itself by injectivity, though possibly not on its residue field. 

These local conditions include all \(H^0\) and inject in \(H^1\) there also. They give exact complementary dualities over the DVR: nullhomotope the cup via its unramified inflation factorization (residue cohomological dimension one with free models). Over a field the two unramified \(H^1\) dimensions equal the respective \(H^0\)'s; by local duality and the local Euler characteristic zero these are orthogonal complements. This and inclusion of all invariants, no \(H^2\), checks the quotient duality equivalence by base change.
\end{itemize}

In all tests global invariants vanish (also dually) by the open-image evaluation argument below. Global Euler characteristic in the problems is \(-1\), by the ordinary formula, or computing the positive-model Euler characteristic \(-2\) stagewise at constant characteristic-zero coefficients and restoring the real term. Unramified conditions at odd places of the third test have Euler correction zero. Consequently the cohomology on the fields has amplitude \(1,2\) by duality and exact orthogonality, with
\[
                         h=1+h^*,
\]
where \(h,h^*\) denote primal and exact-annihilator dual dimensions in degree one. Both degree-one spaces inject into ordinary unrestricted global \(H^1\) (dual strict uses the old dual local invariant vanishing where full old conditions are imposed). These assertions continue to hold with finitely many new line conditions in the argument.

\emph{Local lines at a switching prime.}
Switch at fresh good primes \(\ell=\ell_i\) tending to 1 in \(\mathbb Z_2^\times\), splitting on all base scalar data to increasing precision, with Tate Frobenius limiting to
\[
                \tau=\begin{pmatrix}1&e\\0&1\end{pmatrix},\qquad e\ne0,
\]
deep in an open-image special linear subgroup. Here \(i\) runs over stages of the relevant limit (internal stages in the second and third tests). Inertia acts trivially. On the limit DVR the local complex is split free on cohomology in \(0,1,2\); degree one has axes the finite and singular lines
\[
       f:\ M_{\mathscr O}/(\tau-1)M_{\mathscr O},\qquad
       s:\ \ker(\tau-1\mid M_{\mathscr O}),
\]
by evaluation on a Frobenius lift and a tame inertia generator respectively; \(M_{\mathscr O}\) denotes the coefficient there. Indeed \(e\) is a unit, and one uses the unramified and singular blocks from inertia/residue cochains (the conjugation exponent tends to 1). In particular the pure singular evaluation class sending inertia to the invariant basis and Frobenius to zero can already be computed to growing precisions with \(\ell\equiv1\), Frobenius action \(\equiv\tau\). The two pure lines, taking in addition all \(H^0\) but not \(H^2\), are self-annihilating against their respective opposite versions under the identification by the Weil form (inverse scalar there, also trivial locally in the limit). For the finite line use inflation isotropy; for the pure singular line the just-described cocycles take values along the same Weil-isotropic invariant line, giving the limiting cup zero also. Perfectness then follows by local duality (the mixed pairing is perfect); these are exact orthogonal complex conditions since local cohomology splits, using nullhomotopies as in Section~\ref{sec:foundations}. Use likewise lower taking only degree zero and upper taking both coordinates in addition, exact opposites. The finite condition identifies by inflation with proceeding without allowing this prime. Each intermediate line condition has Euler correction zero and takes in all \(H^0\).

\emph{Selection of the switching primes.} For \(\mathcal G=\prod_i G_\mathbb Q\) take the product \(\mathcal H\) of the exact kernels of Tate action together with metabelian data \(\mathcal E_i\), containing \(\mathbb Q^{\rm ab}\) and the \(2\)-power roots of all old derivative primes (it suffices to take their composite). The scalar actions in use factor through those data. The image quotient contains as a normal subgroup arbitrary sequences of a fixed deep \({\rm SL}_2^b(\mathbb Z_2)\) in the Tate action alone, trivial on \(\mathcal E_i\), by open image for the non-CM curve and closed iterated commutators. Choose \(\tau\) as above in that subgroup (using its constant sequence). Thus by the product-group evaluation and diagonal-vanishing argument in Section~\ref{sec:foundations}, global \(H^0\) vanishes on all fields in question; restrictions to \(\mathcal H\) detect crossed global classes. On \(\mathcal H\) these evaluations are additive, and for a nonzero class the evaluations span the plane by equivariance and absolute irreducibility. These statements apply separately to primal and dual, whose Selmer degree-one classes inject as above.

If on the fiber \(h>1\) with a sole generic line already established, choose the nonzero fiber class from a primitive reduction of that line, and some nonzero degree-one dual class. In a generic-only surplus-rank test below instead take a specified nonzero primal class and some nonzero dual class. In a product-group coset of \(\mathcal H\) having action \(\tau\) and trivial metabelian action, we can simultaneously make both evaluations nonzero modulo \((\tau-1)\): each is a proper affine avoidance test on the additive evaluations, possible jointly by integer combinations in characteristic zero. Coboundary adjustments do not affect the tests. We require Frobenius agreement with the resulting sequence on the old diagram evaluations to growing precisions.

\emph{The depth margin for descent.}
Keep a full finite abelian base \(K_0=\mathbb Q(\mu_{M_i})\) per stage for the Kato constructions with \(2\)-part sufficiently deep for cofinal base scalar precisions, fixed before adding derivative primes. This contains any fixed scalar fields needed for the weighted trace and uses the chosen old omitted support as specified for each system (in square order include the moving primes just once). 

Within the necessary termwise finite open constraints one can still keep \(\mathcal E_i\)-action exactly trivial while perturbing \(\tau\) by arbitrarily deep special linear elements so \(v_2\det(1-\rho_T)\) is large but finite. Indeed intersecting an open subgroup with that metabelian kernel still gives open image inside a deep special linear subgroup; one may perturb in the lower unipotent. \textbf{Then} choose \(A_i\) much larger still. Chebotarev gives \(\ell_i\) with agreements to the requisite finite precisions, splitting in \(K_0\) and in all previously chosen derivative fields, with \(\ell_i\equiv1\bmod 2^{A_i}\), and with
\[
        v_2\det(1-\rho_T({\rm Fr}_{\ell_i}))\longrightarrow\infty,\qquad
        A_i-v_2\det(1-\rho_T({\rm Fr}_{\ell_i}))\longrightarrow\infty .
\]


Impose also splitting on \(2^{A_i}\)-th roots of all previous derivative primes. Thus those primes split in the new degree \(s_\ell=2^{A_i}\) subfield of \(\mathbb Q(\mu_{\ell_i})\) (the new derivative field). Exclude old places, primes of smoothing integers and all level primes already present. This realizes the two nonzero tests by choosing the place representatives (cf. Chebotarev on evaluation matrices in Section~\ref{sec:foundations}). A Frobenius lift trivial also in the new derivative field can be used, by inertia adjustment invisible on the old data needed. 

Further additions keep previous choices and diagram comparisons, reducing to slower cofinal precisions if needed. Splitting in \(K_0\) and the mutual splitting at all other derivative primes here are exact at each stage, not just limiting trace conditions.

\emph{Derivative classes and descent.}
For the square or rectangular Kato system being used, keep its base orders/mark/symbols/smoothings and add each derivative prime once to both orders and the determinant index. For \(I\) a subset of such primes at the current stage let \(K_I/K_0\) be the composite with their derivative fields. Trace first to \(K_I\), with \(T\)-coefficients denoted throughout by their identification there (so transport by the fixed quadratic twist is allowed, trivial as a character over \(K_0\)). Let \(c_I\) denote the integral resulting class, using \eqref{eq:K1}. Write \(\sigma_\ell\) for a tame inertia representative mapping to a generator on its own cyclic derivative factor and trivial on all the other factors. Consider
\[
          D_I^{\rm der} c_I,\qquad
          D_I^{\rm der}=\prod_{\ell\in I} D^{\rm der}_\ell,\qquad
          D^{\rm der}_\ell=\sum_{j=1}^{s_\ell-1}j\sigma_\ell^j .
\]
The good Euler norm of Proposition~\ref{prop:siegel-norms}, pushed by trace to these fields, gives norm down at \(\ell\) with scalar
\(p_\ell=\ell^{-1}\det(1-\rho_T({\rm Fr}_\ell))\); the projection of the field translation is trivial there by the splitting. This is integral equality after a uniform fixed nonzero multiple \(L_0\). Indeed \(H^0((T\otimes\mathbb Q_2)/T)\) over \textbf{all} these abelian fields has bounded exponent by the commutator image, which bounds the torsion ambiguity in \(H^1(T)\); inflation into unrestricted cohomology is injective in degree one. All norm calculations can thus be made with the required enlarged supports also.

Work modulo \(2^n\), \(n\to\infty\) sufficiently slowly that \(s_\ell,\ p_\ell\equiv0\bmod 2^n\) and \(A_i>n\) for all current factors (using each one's chosen depth). The identities \((\sigma_\ell-1)D_\ell^{\rm der}=s_\ell-N_\ell\), \(N_\ell\) the cyclic sum, then show \(L_0 D_I^{\rm der}c_I\) invariant modulo \(2^n\). 

Multiply by a further fixed integer \(L_1\ne0\) for descent by inflation-restriction: the obstruction in degree two of the finite group with coefficients \((T/2^n)^{G_{K_I}}\), and similarly inflation ambiguity, are killed by a bounded power of 2 by the same open-image argument. This applies on unramified-outside-allowed-support groups (including the derivative ramification). Write \(L_*=L_0L_1\), the same for all subsets. 

Descend \(L_*D_I^{\rm der}c_I\) modulo \(2^n\) to \(K_0\), then weighted corestrict from \(K_0\) to the global deformation coefficients (trivialize the scalar over \(K_0\) at the working precisions). Taking limits gives integral ordinary classes \(Z_I\) with \(Z_\varnothing=L_* z_{\rm sys}\), where \(z_{\rm sys}\) is the original system's base class. Here for the empty set we simply use the original representative; compatibility with the true \(t\)-tower uses the \(2\)-power compatibilities of \eqref{eq:K1}. 

Choices need not give an Euler system of descents in \(t\): reduce as needed to cofinal Artin quotients, taking representatives to cycles on the models there and matrix limits. This preserves any previously chosen lower-subset classes through inflation. The weighted corestriction has no division by the group order.

\emph{The finite--singular comparison.}
We give the local comparison, allowing the moduli \(2^n\) to be decreased cofinally. At \(\ell\in I\) take \textbf{integral} cocycles \(b,b_\ell\) on \(G_{K_{I-\{\ell\}}},G_{K_I}\) respectively, for lower and upper classes after applying \(L_*\) and \emph{all other} derivatives. Write \(\sigma=\sigma_\ell,\ m=s_\ell,\ \Phi={\rm Fr}_\ell\) a lift trivial on \(K_I\), and \(F=\rho_T(\Phi)\). 

Both cocycles vanish on the respective local inertias: inertia acts trivially on \(T\), and even over the upper cyclic ramification \(H^1(I_{\rm loc},T)^\Phi=0\) by good reduction and Weil (no eigenvalue \(\ell\); the residue degree here is one). This applies at conjugate places as well. Put \(x=b(\Phi), x_\ell=b_\ell(\Phi)\). On \(G_{K_I}\) the exact norm equality gives, for some \(y\in T\),
\[
 \sum_{j=0}^{m-1}(\sigma^j b_\ell)(u)=p_\ell b(u)+(u-1)y,\qquad
 (F-1)y=m x_\ell-p_\ell x.
\]
Here a translate conjugates the argument and acts on the coefficient. All the translates in the sum have the same value on \(\Phi\), by the inertia vanishing and trivial inertia action. Also \(\sigma^m b_\ell=b_\ell\) as cocycles by the same vanishing on \(\sigma^m\).

Let \(B_I\) be the mod-\(2^n\) descended cocycle before weighted corestriction. On the upper subgroup adjust by a coboundary (adjusting also on the whole group of \(K_0\)) to make its restriction equal to \(D_\ell^{\rm der}b_\ell\). Its value on \(\Phi\) after adjustment is zero by \(m(m-1)/2\equiv0\). Applying \(\sigma-1\) on these restrictions gives, for \(u\in G_{K_I}\),
\[
                 (u-1)B_I(\sigma)=-(u-1)y\bmod 2^n,
\]
using the crossed identity, the norm calculation and \(\sigma^m b_\ell=b_\ell\). Thus \(B_I(\sigma)\equiv-y\) losing at most bounded \(2\)-power precision by the invariant bound, independently of that coboundary adjustment. Now
\[
             y=m(F-1)^{-1}x_\ell-\ell^{-1}{\rm adj}(F-1)x .
\]
The first term vanishes in the limit by the margin on the determinant valuation. And \(x\) agrees, up to a coboundary coordinate \((F-1)w\) at the common precision, with \(B_{I-\{\ell\}}(\Phi)\). Such changes disappear after applying \({\rm adj}(F-1)\) in the limit.

Upon restriction of weighted corestriction to \(G_{K_0}\), one obtains the sum of translates with weights. On \textbf{each translate} the same calculation applies with the same local representatives: the translated lower/upper classes still have the derivative and norm identities in cohomology (all factor actions commute there by abelianness), and inertia vanishing holds. Hence after summation and limit
\[
         f(Z_I)=0,\qquad s(Z_I)={\rm adj}(\tau-1) f(Z_{I-\{\ell\}})
                                    \quad\text{at }\ell .           \tag{L3}\label{eq:L3}
\]
All congruences here concern singular evaluations with trivial coefficient inertia and finite evaluations modulo Frobenius coboundaries. They therefore carry to the cohomology maps on the limiting models, then to the DVR or test field. The verification can likewise be made at each old \(\ell\in I\) separately to keep the pure singular conditions there. 

The line map in \eqref{eq:L3} is a unit isomorphism over \(\mathscr O\) (off-diagonal \(\pm e\)). The \(Z_I\) lift in cohomology over the DVR with the prescribed pure conditions: degree zero is included and degree-one coordinates have the prescribed image. In the third test their old odd-place lifts to unramified conditions before residue-field specialization are automatic by zero \(H^1\) of the singular quotient over the DVR. The same facts apply using just generic fields.

\emph{Propagation of determinant divisibility.}
If in a generic test \(z_{\rm sys}\ne0\), there is no surplus \(h>1\) there: evaluate it and an old nonzero dual class as above at a first fresh prime; by \eqref{eq:L3} the derivative is then nonzero transverse, contradicting local cup reciprocity with that dual class (whose finite localization there is nonzero). 

In a DVR test with the single generic line and nonzero class established, proceed successively if \(h>1\) on the fiber. At each new prime start with unramified inflation, arrange the two evaluations as above on the fiber (so also nonzero generic primal evaluation), and switch by the \textbf{Line switch} comparison. This drops \(h\) by one preserving generic nonvanishing of the class by \eqref{eq:L3}. The ordinary determinant valuations are exactly equal along the comparisons with the new derivatives; there is no bound lost at this DVR. 

When \(h=1\) the terminal complex is a free line in degree one up to homotopy and the terminal class is integral. Since \(L_*\) is a unit here, backwards we obtain
\[
                         d(C_{\rm old},z_{\rm sys})\ \ge\ 0.         \tag{L4}\label{eq:L4}
\]
Here \(C_{\rm old}\) is the initial ordinary problem in the switch sequence,
with the degree-one lift when unramified conditions are imposed. The clearing
factor \(L_*\) is a fixed power of two up to a unit, and is a unit at each
DVR used here. Consequently the equality of determinant valuations in the
switches gives exactly \eqref{eq:L4}.
\end{proof}

\begin{corollary}[Generic nonvanishing and the first divisibilities]\label{cor:positive-determinant:generic}
The positive complex has only degree-one generic cohomology, of dimension
two, and \(a\wedge z\ne0\). The coordinate \(U\) has nonnegative order at
every height-one prime outside the fixed factors
\(2\Delta\prod_{q\in S_f,\,q\ne2}D_qP_q\).
\end{corollary}
\begin{proof}

Apply first the fixed-support \(t\)-test with square order and \(z_{\rm sys}=z\) there. Its ordinary class is generically nonzero: otherwise a nonzero bounded series annihilates it and it specializes to zero rationally off finitely many finite-order characters. This contradicts \eqref{eq:K4} and cyclotomic twist nonvanishing~\cite{Rohrlich88}, in the form of \cite[Theorem~13.5(2)]{Kato04}, at high characters, including with the fixed quadratic twist when in use. Thus the test gives \(h=1\) generically. 

Now specialize the \(B\)-diagram at \(u_j=0\) for all \(j\), retaining generic \(t\). The moving places have unramified action there and acyclic singular quotient (cyclotomic Frobenius scalar tends to 1, residue sizes tend to 1 and \(\det(1-g_j)\ne0\)). Inflation thus compares isomorphically with the fixed-support problem. Moreover the augmented ordinary square class is that from fixed support times nonzero limiting Euler factors: apply the fresh good norms removing the \(r_{j,i}\) sequentially at trivial tame weight, with bounded torsion clearing as above; \(\chi_h\) if used is split there, and the factors tend to \(2-\operatorname{tr}g_j\). 

Hence there is a single ordinary line with nonzero class on this specialized generic line. By deformation, or using also the positive concentration at \(B_{(2)}\), we have generically on full \(B\) just the two-dimensional degree one in \(C^+\), with \(a\wedge z\ne0\) (the specialized generic \(t\)-test has these same dimensions). This holds also for analytically unknown center. 

The first, \(B\)-DVR test now gives \eqref{eq:L4} for the ordinary full complex there, hence nonnegative order for \eqref{eq:L1} away from \((2)\) and the indicated fixed factors.
\end{proof}

\subsection{Removal of the fixed divisors}\label{subsec:positive-determinant:division}

\begin{lemma}[Division at the fixed factors]\label{lem:positive-determinant:division}
The coordinate \(U\) has nonnegative order at every horizontal height-one
prime dividing \(\Delta\prod_{q\in S_f,\,q\ne2}D_qP_q\).
\end{lemma}
\begin{proof}
\emph{Finite-character tests.}

Represent \(U=g_0/f_0\) with integral series \(f_0,g_0\), \(f_0\ne0\), by determinant pivot formulas from the positive diagram, incorporating the factors of \eqref{eq:L1}. Retain their stage expressions \(g_{0,i},f_{0,i}\) before passage from the actual \(t\)-towers with finite tame orders. These use the integral smoothed lifts and basis choices; they compute the indicated rational coordinate when the pivots work with the same ranks. 

Test at \(u_j=\eta_j-1\) for finite \(2\)-power roots \(\eta_j\ne1\), requiring \(f_0(\cdot,\boldsymbol\eta-1)\ne0\) and dual nonexceptionality on the Frobenius lift at 2 as in the third switching test. The latter is a nonidentical vanishing avoidance condition by the nonzero limiting tame-scalar exponent vector there. These tests still detect nonzero bounded series conditions on the tame variables by iterated one-variable Weierstrass. At each such test we have all the required finite characters and 2-local constant-scalar agreement at sufficiently precise original stages on the filter.

\emph{A rectangular system with the primitive odd factors.}
Fix one such stage for the third test. To retain the precise primitive odd factors use \eqref{eq:K1}, \eqref{eq:K4} in the \textbf{rectangular} first-row-unfilled case. Take \(b=N_{\rm odd}\) and \(M\) supported only at 2 and the odd ramified places of the constant field weight (all active variable conductors and any active fixed quadratic ramification). These are coprime to \(b\); use sufficiently high \(2\)-part and a true \(2\)-power tower. Choose \(A=2^k,\xi\) with \(\delta_\xi(a)\ne0\) as in Lemma~\ref{lem:siegel-rectangular-symbol}. 

At any given height-one \(\mathfrak q\) of \(\Lambda_\eta\) not over 2, take smoothing \(c',d'\) whose factors \(\Delta'\) are units there. Indeed \(t\) at such a divisor can be realized at an algebraic point of the open unit disk by preparation. Take a CRT sequence tending dyadically to 1, congruent to 1 on required fixed symbol orders, respecting coprimalities and realizing a fixed nontrivial constant weight on an active variable conductor (values trivial on all others and any fixed quadratic factor). This uses the disjointness of the active variable conductors from those fixed levels. The \(t\)-scalar tends to 1; thus the vanishing equation in the smoothing is avoided. 

Keep the choices fixed while switching. Write \(z'\) for this system. It has no extra odd imprimitive losses at the base finite cyclotomic characters: \eqref{eq:K4} gives the law with \(L^{(2)}\) since all omitted odd places are at ramified scalar characters of good-\(f\) primes. In particular its generic class is nonzero by the same cyclotomic nonvanishing theorem. The rectangular good Euler norm is exactly of the requisite kind. Thus generic single-line concentration for the unramified-old-odd problem, and \eqref{eq:L4} with \(z'\) at \(\mathfrak q\), follow by the third switching test. The same generic concentration holds on relaxing to full, since old odd singular quotient differentials have nonzero determinants.

\emph{Comparison of the two classes.}
In this common generic full cohomology line the original square system differs from \(z'\) by
\[
             \frac{\Delta\,\prod_{q\in S_f,\ q\ne2}P_q}{\Delta'}
\]
times a nonzero constant, using the stage specializations for these factors. Indeed evaluate at infinitely many high \(t\)-characters with central nonvanishing, using \eqref{eq:K4} in the \textbf{same} representation with compatible weighted realizations; any geometric quadratic identification is common. The two exponential coordinates have precisely the displayed ratio times the fixed symbol ratio, up to signs. Odd omissions at the active variable conductors make no change. Generic ratios specialize in ordinary cohomology off at most finitely many such values, by clearing denominators, so the identity follows by one-variable Weierstrass (one can square to ignore signs). In particular the square raw coordinate uses a single generic ordinary line also on the tested stages. 

On relaxation from the unramified-old-odd complex to full, determinant valuation changes for \(z'\) by the sum of \(d(Q_q,1)=-v_{\mathfrak q}(D_{q,\eta})\) for those singular quotients. These determinants are 1 at active variable conductors, and are precisely the specialized stage \(D_q\)'s at fixed odd places. 

Consequently \eqref{eq:L4}, the ratio comparison, and the ordinary real basis give
\[
 \frac{g_{0,i}(\cdot,\boldsymbol\eta-1)}
      {f_{0,i}(\cdot,\boldsymbol\eta-1)}
                 \ \in\ \mathbb Z_2[\eta][[t]][1/2].
\]
Here we use all the horizontal height-one primes, allowing \(z'\) or its smoothing choices to vary among them. On large stages the chosen pivots indeed have nonzero denominator generically on the test, by coefficientwise convergence and \(f_0(\cdot,\boldsymbol\eta-1)\ne0\), and ranks agree as just proved. No bound on a comparison constant or uniform \(2\)-denominator is being asserted, and this step uses no separate primitive-form characteristic-ideal divisibility theorem.

\emph{Integral Weierstrass division.}
The remaining possible prime divisors of poles divide \(\Delta,D_q,P_q\) at fixed odd corrections (or the smoothing alone); every such factor is regular in \(t\) over \(\mathbb Z_2[[\mathbf u]]\) as observed after \eqref{eq:L1}. 

For each irreducible such factor \(H_0\) use its distinguished representative, and \(H_0^a\mid f_0\) with \(a=v_{H_0}(f_0)\). At the above tests it stays distinguished. All its tested zeros to that multiplicity belong to \(g_0(\cdot,\boldsymbol\eta-1)\), by the integral numerator/denominator zero-persistence test in Section~\ref{sec:foundations} (convergence on smaller disks and the stage containment). 

Thus the integral division remainder of \(g_0\) by \(H_0^a\) in \(t\) vanishes at every such test, hence identically. This rules out the remaining horizontal poles.
\end{proof}

\begin{proof}[Proof of Proposition~\ref{prop:positive-determinant:integral}]
Lemma~\ref{lem:positive-determinant:residual} proves concentration and
nonnegative order at \((2)\). Corollary~\ref{cor:positive-determinant:generic}
proves generic nonvanishing and the required divisibility away from the fixed
factors, and Lemma~\ref{lem:positive-determinant:division} treats those factors.
The ring \(B\) is normal, so these height-one tests give \(U\in B\).
Generic nonvanishing gives \(U\ne0\). Finally, the common raw reduction of
Lemma~\ref{lem:positive-determinant:residual}, the common smoothing reduction,
and Lemma~\ref{lem:positive-determinant:factors} show that the two integral
reductions differ by a unit. They therefore have the same closed-point unit
status, as asserted.
\end{proof}

\subsection{The central value in Selmer corank zero}\label{subsec:positive-determinant:rank-zero}

\begin{lemma}[Corank-zero central specialization]\label{lem:positive-determinant:rank-zero}
For either branch \(A_0\) of the construction, if \(\mathrm{an}(A_0)=0\),
then the normalized determinant satisfies
\[
                         v_2 U(0)=X(A_0) .                           \tag{L5}\label{eq:L5}
\]
In particular \(X(A_0)\ge0\).
If only \(s_2(A_0)=0\) is known and \(U\) is a unit, then
\(L(A_0,1)\ne0\), and hence \(X(A_0)=0\).
\end{lemma}
\begin{proof}
\emph{Compact and discrete Selmer lattices.}

For the central calculation first work at the actual trivial-parameter specialization at stages, for the curve of the branch, denoted temporarily by \(A_0\). Put \(S'=S\cup R_i\), \(O^j=H^j(G_{\mathbb Q,S'},T)\), \(R^j_\mathbb R=H^j(\mathbb R,T)\). Suppose \(s_2(A_0)=0\), which in particular holds for analytic rank zero. 

Put \(s_0=\operatorname{length}\Sha(A_0/\mathbb Q)[2^\infty]<\infty\), and write \(\tau_g,\tau_v\) for lengths of global and finite local \(2\)-primary point torsion. Lengths here are over \(\mathbb Z_2\). The usual discrete Selmer has length \(s_0\). At compact coefficients the singular quotient by Kummer of local degree one is zero at odd and real places, and at 2 it is \(J_2\) free of rank one, dual to compact Kummer modulo torsion. 

With \(j_2=\operatorname{length}(J_2/\operatorname{im}O^1)\), integral Poitou--Tate gives
\[
       0\le j_2\le s_0,\qquad
       \operatorname{length}O^2
             =s_0-j_2+\operatorname{length} R^2_\mathbb R
                        +\sum_{v\in S'_f}\tau_v-\tau_g .
\]
Indeed the exact sequence goes from \(O^1\) to the local singular quotient, the dual of discrete Selmer, \(O^2\), the sum of local degree-two groups and the dual of global point torsion, ending surjectively. This uses modified real localization in Poitou--Tate (agreeing with ordinary real terms in those positive degrees), finite Kummer duality and passage to compact/discrete limits. At odd good places outside support unramified discrete \(H^1\) is already zero, giving the usual Selmer. 

Also \(O^0=0,\ \operatorname{rank}O^1=1,\ \operatorname{length}O^1_{\rm tors}=\tau_g\), by the Euler and invariant formulas.

\emph{The determinant and local volume factors.}
Assume first that the analytic rank is zero. Then \(z\) has nonzero singular localization by \eqref{eq:K4}; write \(z\) also for its ordinary image. The positive complex with the integral real boundary basis and lift then gives raw valuation
\begin{equation}\label{eq:positive-determinant:raw-center}
 \begin{aligned}
 d(C^+,a\wedge z)
   &=\operatorname{length}(O^1/\mathbb Z_2 z)
       -\operatorname{length}O^2-\operatorname{length}R^1_\mathbb R
                                  +\operatorname{length}R^2_\mathbb R\\
   &=v_2(z,J_2)+2\tau_g-s_0-\sum_{v\in S'_f}\tau_v
                                       -\operatorname{length}R^1_\mathbb R ,
 \end{aligned}
\end{equation}
where \(v_2(z,J_2)\) measures the singular image in that lattice. The first equality uses the long exact positive sequence through degree three: dividing degree one by the image of the real basis gives the kernel to \(R^1_\mathbb R\), and positive degree three is the cokernel to \(R^2_\mathbb R\). The second uses \(\operatorname{length}(O^1/\mathbb Z_2z)=\tau_g+v_2(z,J_2)-j_2\).

By the differential/log duality and Haar computation of the central lattices, for minimal \(\omega\), \(\exp^*z=\alpha\omega,\ \log_\omega A_0(\mathbb Q_2)=2^b\mathbb Z_2\), we have \(v_2(z,J_2)=v_2(\alpha)+b\), and
\[
       \tau_2-b=v_2(c_2 L_2(A_0,1)^{-1}),\qquad
       \tau_q=v_2(c_q L_q(A_0,1)^{-1})\quad(q\ne2).
\]
The \(R^1_\mathbb R\) length inserts the real component factor. Thus the trivial-parameter law of \eqref{eq:K4} on the indicated curve (\(\Delta(0)\), omitted support \(S'_f\), connected period, evaluation 1) gives exactly \(X(A_0)+v_2\Delta(0)\) for the raw valuation. 

\emph{Specialization of the limit determinant.}
These ranks and the determinant calculation pass to the limit at the center. Indeed the moving finite torsion lengths stay bounded by the trace-limit conditions. The formulas above bound \(O^2\), hence also positive degrees two and three; torsion in positive degree one injects into ordinary torsion. Pivots of the same ranks therefore work on the limiting specialization, allowing comparison of determinant and cycle entries there by bounded-denominator Gaussian formulas. Thus the raw valuation is unchanged in analytic rank zero, and the fixed \(D_q/P_q\) have central ratio 1, proving \eqref{eq:L5}.

\emph{An initially unknown analytic center.}
In particular \(X(A_0)\ge0\) in analytic rank zero. If only \(s_2(A_0)=0\) is known but \(U\) is a unit, the same torsion bounds and ranks apply regardless of analytic nonvanishing. If \(L(A_0,1)=0\), then \(z\) has trivial rational singular image by \eqref{eq:K4}, hence is ordinary torsion at the stages, so \(a\wedge z=0\) rationally there and also at the limiting center, a contradiction (all fixed corrections and smoothing have nonzero central values). Thus \(L(A_0,1)\ne0\); now \eqref{eq:L5} applies and gives \(X(A_0)=0\).
\end{proof}

In rank one, computing the same determinant requires a comparison of the Kato class with a regulator generator. Section~\ref{sec:tame} proves that comparison, and Section~\ref{sec:positive-completion} applies it.

